\section{Extensional solution. Lateral forces.} % 348
In a conical shell, when
\(n=1\) and
\(g_1\), \(g_2\), \(h_1\)
are neglected, equations (4) become
\begin{equation}\label{eq92}
\begin{array}{l}{{\displaystyle\frac{d}{d r}\left(t_{1}r\right)\sin\gamma+s_{1}
-t_{2}\sin\gamma+\frac{1-\sigma^{2}}{2E h}\,r X^{\prime\prime}\sin\gamma=0,}}\\
{{\displaystyle\frac{d}{d r}\left(s_{1}r\right)\sin\gamma-t_{2}+s_{1}\sin\gamma+
\frac{1-\sigma^{2}}{2E h}\,r Y^{\prime\prime}\sin\gamma=0,}}\\
{{\displaystyle
\frac{d}{d r}\left(s_{1}r\right)\sin\gamma
t_{2} +\frac{1-\sigma^{2}}{2E h}\,r Z^{\prime\prime}\tan\gamma=0.}}\end{array}
\end{equation}

On multiplying the 1st of equations~(\ref{eq92}) by
\(\sin \gamma \), the 2nd by
\(-1\), and the 3rd by
\(\cos^2 \gamma \), and adding, we obtain the equation
\[
{\frac{d}{d r}}(t_{1}r)\ \sin^{2}\gamma-{\frac{d}{d r}}(s_{1}r)\sin\gamma+{\frac{1-\sigma^{2}}{2E\hbar^{2}}}\,r\,(X^{\prime\prime}\sin\gamma-Y^{\prime\prime}-Z^{\prime\prime}\cos\gamma)\sin\gamma=0,
\]
which can be integrated in the form
\begin{equation}\label{eq93}
t_{1}r\sin\gamma-s_{1}r=a_{1}-{\frac{1-\sigma^{*}}{2E\hbar}}\int_{r_{\omega}}^{r}r\left(X^{\prime\prime}\sin\gamma-Y^{\prime\prime}-Z^{\prime\prime}\cos\gamma \right)d r
\end{equation}

The elimination of \(t_2\) between the 1st and 2nd of equations~(\ref{eq92}) gives
\[
{\frac{d}{d r}}\,(t_{1}r)\sin\gamma-{\frac{d}{d r}}\,(s_{1}r)\sin^{s}\gamma+s_{1}\cos^{s}\gamma+{\frac{1-\sigma^{s}}{2E{\hbar}}}\,{}^{s}(X^{\prime\prime}-Y^{\prime\prime}\sin\gamma)\sin\gamma=0,
\]
and the elimination of \(t_1\) between this equation and equation~(\ref{eq93}) gives
\[
\left(\frac{d}{d r}\left(s_{1}r\right)+s_{1}\right)
\cos^{2}\gamma+\frac{1-\sigma^{2}}{2\bar{E}\hbar}\,r\left(Y^{\prime\prime}\cos^{2}\gamma+Z^{\prime\prime}\cos\gamma \right)=0,
\]
which can be integrated in the form
\begin{equation}\label{eq94}
s_{1}r^{2}\cos\gamma=a_{2}-\mathrm{\frac{~1-\sigma^{2}}{~2E h}}\int_{r_{0}}^{r}r^{2}\left(Y^{\prime\prime}\cos\gamma+Z^{\prime\prime}\right)d r.
\end{equation}

From equations equations~(\ref{eq93}) and~(\ref{eq94})
\(t_1\) and
\(s_1\) are known, and
\(t_2\) is known from the 3rd of equations~(\ref{eq92}).
\todo[inline]{fixed equation part 3 again}

Now equations (2) and (3) give
\begin{equation}\label{eq95}
\begin{array}{c}
{{ t_{1}= {\displaystyle\frac{d U}{d r}}
+\sigma\left({\displaystyle\frac{U}{r}}+{\displaystyle\cfrac{V}{r\sin\gamma}}-{\displaystyle\frac{{\bar{v}}}{r\tan\gamma}}\right), }} \\
{ {\displaystyle t_{2}= \sigma{\displaystyle\frac{d U}{d r}}+
\left({\displaystyle\frac{U}{r}}+{\displaystyle\cfrac{V}{r\sin\gamma}}-{\displaystyle\cfrac{U}{r\tan\gamma}}\right).} } \\
{ 
{\displaystyle s_{2}= \frac{1-\sigma}{2} 
\left( \sigma {\displaystyle\frac{d V}{d r}}-
{\displaystyle\frac{V}{r}}+{\displaystyle\cfrac{U}{r\sin\gamma}} 
\right)
} }
\end{array}
\end{equation}

From these we find
\[ {\frac{d U}{d r}}={\frac{t_{1}-\sigma t_{2}}{1-\sigma^{2}}}\,, \]
whence
\begin{equation}\label{eq96}
U=b_{1}+\int_{r_{0}}^{r}\frac{t_{1}-\sigma t_{2}}{1-\sigma^{2}}\,d r\,;
\end{equation}
